\section{Where gsplit Underperforms}
\label{sec:limitations}

\paraq{What this section establishes, and why it matters.}
The starting point is not that gsplit is a poor design: in our earlier results it was generally the
fastest option across the grid, and we treat that as a prior to re-examine here, not a given. What
this section maps is how gsplit's \emph{advantage} over the DGL baseline behaves as we vary the
interconnect, the graph, the model, and the model's width and depth. The motivating observation for
the rest of the thesis is that this advantage is not uniform---in parts of the space it
\emph{shrinks}, and in some it \emph{vanishes or reverses}. An advantage that is absent, or present
but diminishing, is already enough to motivate the hybrid-parallelism search of
Section~\ref{sec:hybrids}: wherever a single fixed strategy stops pulling ahead, a configuration that
mixes strategies may recover the headroom. We therefore characterise the space before drawing a
conclusion from it.

\paraq{The comparison, and how caching is handled.}
We use a single head-to-head: \textbf{gsplit} (\(k{=}0\)) against \textbf{DGL}, both at a $0\%$
feature cache, so that only the parallelization strategy differs. A feature cache is a large,
system-specific effect that would otherwise confound the comparison; rather than pit gsplit against a
cache-based baseline, we hold the cache at zero for both and report, alongside the per-epoch time, the
\emph{training-only} time (forward${}+{}$backward) in parentheses. The training-only figure is the
component that would survive a perfect feature cache---if caching made feature loading free, the
per-epoch time would collapse toward it---so the pair of numbers brackets gsplit's behaviour with and
without the data-movement cost that a cache removes. We read $\Delta$ as gsplit relative to DGL,
\textbf{positive $\Rightarrow$ gsplit is faster} (a gain); a $\Delta$ at or below zero says the gain
is gone, and a positive $\Delta$ that \emph{shrinks} along a ladder says the gain is diminishing.
Both outcomes motivate Section~\ref{sec:hybrids}.

\paraq{Experimental setup.}
All measurements run on a single node of $4{\times}$\textbf{A100-SXM-80GB} GPUs. NVLink is the
native interconnect (NV12, all pairs); the PCIe-class setting is \emph{emulated} on the same node by
disabling NCCL peer-to-peer (\texttt{NCCL\_P2P\_DISABLE=1}), which forces the same shuffles onto the
slower path without a hardware change. The two systems are \textbf{DGL} (full-graph replication,
features gathered over UVA from host-pinned memory) and \textbf{gsplit} (\(k{=}0\)), both with
\emph{no} feature cache. Training uses the \texttt{sage} and \texttt{gat} models, neighbour fanout
\textbf{15} per layer, batch size \textbf{1024}, and $\text{world size}{=}4$. Because we measure
\emph{time}, not accuracy, features are synthetic (\texttt{--random\_feat}: a full-size tensor, so
feature movement stays realistic while the real feature load is skipped). Each cell runs
\textbf{3 epochs} and we report the mean \emph{per-epoch} time (\texttt{duration (s)}), with the
training-only time (forward${}+{}$backward) in parentheses; a per-cell wall-clock cap of
\textbf{1200\,s} records a \textsc{t/o}, and a cell that exhausts GPU memory records \textsc{oom}. We
stress both \textbf{products} and \textbf{papers100M}; arxiv is omitted as a correctness-only toy
graph.

\paraq{The measurement matrix.}
Each graph gets one table (Tables~\ref{tab:sec1-products}--\ref{tab:sec1-papers}), with \textbf{PCIe}
(emulated) and \textbf{NVLink} as column groups and \textbf{GAT} and \textbf{SAGE} as the two
horizontal sections. The two stress axes are ablated \emph{separately}: a \textbf{hidden-size} ladder
$h\in\{256,1024,4096\}$ at depth $L{=}4$, and a \textbf{depth} ladder $L\in\{2,4,6\}$ at hidden
$h{=}1024$, meeting at the shared pivot $(L{=}4,\,h{=}1024)$. Both axes scale either the volume a
shuffle moves or the number of shuffles, so they are the natural places to expect gsplit's advantage
to erode. Both tables are generated from the single consolidated results file
(\texttt{results/sweep\_all.csv}) by \texttt{gen\_sec1\_tables.py}, so the ablation numbers and any
other use of the same runs cannot drift apart. \expected{on PCIe, gsplit's advantage is largest at
small $h$/shallow $L$ and shrinks as either grows; on NVLink the advantage is smaller throughout and
vanishes or reverses in the large-$h$/deep corners.}

% AUTO-GENERATED by defense_latex/gen_sec1_tables.py -- do not edit by hand.
\begin{table}[t]
\centering
\caption{\textbf{products.} Per-epoch time (s); \emph{training-only} time
(forward+backward) in parentheses --- the compute that survives a perfect feature
cache. $\Delta$ is gsplit's improvement over DGL (\textbf{positive $=$ gsplit faster}),
shown as per-epoch $(\text{training})$. Two single-axis ablations per model crossing at
the \textbf{default} $(L{=}4,\,h{=}1024)$ --- the value \emph{spanned} (multirow) in each
block: a hidden sweep at $L{=}4$, then a depth sweep at $h{=}1024$. gsplit runs at a
$0\%$ feature cache; PCIe is emulated via \texttt{NCCL\_P2P\_DISABLE=1}. Features random.}
\label{tab:sec1-products}
\setlength{\tabcolsep}{4pt}
\resizebox{\textwidth}{!}{%
\begin{tabular}{|l|c|c|r|r|r|r|r|r|}
\hline
\multicolumn{3}{|c|}{} & \multicolumn{3}{c|}{\textbf{PCIe (emul.)}} & \multicolumn{3}{c|}{\textbf{NVLink}} \\
\cline{4-9}
\textbf{Model} & $\mathbf{h}$ & $\mathbf{L}$ & DGL & gsplit & $\Delta$ & DGL & gsplit & $\Delta$ \\
\hline
\multirow{6}{*}{\textbf{SAGE}} & $256$ & \multirow{3}{*}{$4$} & 13.19 (4.84) & 15.53 (9.51) & $-18\%\,(-96\%)$ & 13.17 (4.80) & 10.97 (4.72) & $+17\%\,(+2\%)$ \\
 & $1024$ &  & 26.08 (18.14) & 38.77 (32.61) & $-49\%\,(-80\%)$ & 25.98 (17.93) & 23.03 (16.59) & $+11\%\,(+7\%)$ \\
 & $4096$ &  & 162.34 (153.97) & 202.71 (196.10) & $-25\%\,(-27\%)$ & 163.36 (155.05) & 147.24 (140.41) & $+10\%\,(+9\%)$ \\
\cline{2-9}
 & \multirow{3}{*}{$1024$} & $2$ & 1.38 (0.64) & 1.76 (0.90) & $-28\%\,(-41\%)$ & 1.34 (0.62) & 1.99 (0.80) & $-49\%\,(-29\%)$ \\
 &  & $4$ & 26.08 (18.14) & 38.77 (32.61) & $-49\%\,(-80\%)$ & 25.98 (17.93) & 23.03 (16.59) & $+11\%\,(+7\%)$ \\
 &  & $6$ & 363.66 (331.17) & 260.82 (246.70) & $+28\%\,(+26\%)$ & 364.76 (332.42) & 177.49 (163.33) & $+51\%\,(+51\%)$ \\
\hline
\multirow{6}{*}{\textbf{GAT}} & $256$ & \multirow{3}{*}{$4$} & 23.34 (15.09) & 21.97 (15.98) & $+6\%\,(-6\%)$ & 23.38 (15.11) & 17.46 (11.25) & $+25\%\,(+26\%)$ \\
 & $1024$ &  & 79.50 (71.62) & 70.85 (64.72) & $+11\%\,(+10\%)$ & 79.48 (71.62) & 55.33 (48.96) & $+30\%\,(+32\%)$ \\
 & $4096$ &  & \textsc{OOM} & \textsc{t/o} & -- & \textsc{OOM} & \textsc{t/o} & -- \\
\cline{2-9}
 & \multirow{3}{*}{$1024$} & $2$ & 2.56 (1.95) & 3.38 (2.49) & $-32\%\,(-28\%)$ & 2.55 (1.93) & 3.18 (2.04) & $-25\%\,(-6\%)$ \\
 &  & $4$ & 79.50 (71.62) & 70.85 (64.72) & $+11\%\,(+10\%)$ & 79.48 (71.62) & 55.33 (48.96) & $+30\%\,(+32\%)$ \\
 &  & $6$ & \textsc{OOM} & 320.56 (306.43) & -- & \textsc{OOM} & 237.36 (223.23) & -- \\
\hline
\end{tabular}%
}
\end{table}

% AUTO-GENERATED by defense_latex/gen_sec1_tables.py -- do not edit by hand.
\begin{table}[t]
\centering
\caption{\textbf{papers100M.} Per-epoch time (s); \emph{training-only} time
(forward+backward) in parentheses --- the compute that survives a perfect feature
cache. $\Delta$ is gsplit's improvement over DGL (\textbf{positive $=$ gsplit faster}),
shown as per-epoch $(\text{training})$. Two single-axis ablations per model crossing at
the \textbf{default} $(L{=}3,\,h{=}1024)$ --- the value \emph{spanned} (multirow) in each
block: a hidden sweep at $L{=}3$, then a depth sweep at $h{=}1024$. gsplit runs at a
$0\%$ feature cache; PCIe is emulated via \texttt{NCCL\_P2P\_DISABLE=1}. Features random.}
\label{tab:sec1-papers}
\setlength{\tabcolsep}{4pt}
\resizebox{\textwidth}{!}{%
\begin{tabular}{|l|c|c|r|r|r|r|r|r|}
\hline
\multicolumn{3}{|c|}{} & \multicolumn{3}{c|}{\textbf{PCIe (emul.)}} & \multicolumn{3}{c|}{\textbf{NVLink}} \\
\cline{4-9}
\textbf{Model} & $\mathbf{h}$ & $\mathbf{L}$ & DGL & gsplit & $\Delta$ & DGL & gsplit & $\Delta$ \\
\hline
\multirow{6}{*}{\textbf{SAGE}} & $256$ & \multirow{3}{*}{$3$} & 12.47 (5.41) & 13.86 (6.59) & $-11\%\,(-22\%)$ & 11.66 (5.05) & 13.36 (5.76) & $-15\%\,(-14\%)$ \\
 & $1024$ &  & 14.96 (8.27) & 21.37 (14.69) & $-43\%\,(-78\%)$ & 14.23 (7.46) & 16.90 (9.34) & $-19\%\,(-25\%)$ \\
 & $4096$ &  & 67.54 (60.34) & 78.83 (70.24) & $-17\%\,(-16\%)$ & 61.22 (53.98) & 62.59 (53.21) & $-2\%\,(+1\%)$ \\
\cline{2-9}
 & \multirow{3}{*}{$1024$} & $2$ & 6.80 (3.34) & 8.40 (4.60) & $-24\%\,(-38\%)$ & 6.72 (3.10) & 7.89 (4.29) & $-17\%\,(-38\%)$ \\
 &  & $3$ & 14.96 (8.27) & 21.37 (14.69) & $-43\%\,(-78\%)$ & 14.23 (7.46) & 16.90 (9.34) & $-19\%\,(-25\%)$ \\
 &  & $4$ & 59.78 (39.95) & 76.93 (59.20) & $-29\%\,(-48\%)$ & 57.93 (38.24) & 51.12 (34.66) & $+12\%\,(+9\%)$ \\
\hline
\multirow{6}{*}{\textbf{GAT}} & $256$ & \multirow{3}{*}{$3$} & 19.69 (13.03) & 22.05 (14.77) & $-12\%\,(-13\%)$ & 19.39 (12.61) & 21.21 (13.52) & $-9\%\,(-7\%)$ \\
 & $1024$ &  & 35.91 (28.84) & 34.89 (28.24) & $+3\%\,(+2\%)$ & 35.16 (28.13) & 30.24 (23.15) & $+14\%\,(+18\%)$ \\
 & $4096$ &  & 197.87 (190.38) & 166.23 (158.05) & $+16\%\,(+17\%)$ & 197.67 (190.24) & 147.96 (139.93) & $+25\%\,(+26\%)$ \\
\cline{2-9}
 & \multirow{3}{*}{$1024$} & $2$ & 11.30 (7.92) & 12.66 (8.83) & $-12\%\,(-11\%)$ & 11.33 (7.75) & 13.27 (9.08) & $-17\%\,(-17\%)$ \\
 &  & $3$ & 35.91 (28.84) & 34.89 (28.24) & $+3\%\,(+2\%)$ & 35.16 (28.13) & 30.24 (23.15) & $+14\%\,(+18\%)$ \\
 &  & $4$ & 135.80 (115.33) & 120.49 (104.27) & $+11\%\,(+10\%)$ & 134.65 (114.33) & 101.44 (83.82) & $+25\%\,(+27\%)$ \\
\hline
\end{tabular}%
}
\end{table}


\paraq{Reading the tables (observations, not conclusions).}
These are readings of the matrix, offered with an explicit element of doubt: the per-epoch times carry
a few percent of noise (warm-up epoch, allocator effects), and GNN performance is strongly
\emph{graph-dependent}, so a shape seen on one graph/interconnect/model is not assumed to hold on
another. Read $\Delta$ \emph{within} a column and along a single ladder: a positive $\Delta$ that
falls as $h$ or $L$ grows is a \emph{diminishing} gain; a $\Delta$ that reaches zero or turns negative
is a cell where gsplit no longer helps. The per-epoch and training-only $\Delta$s are reported
together so that a gain which exists only because of data movement (and would disappear under a cache)
is visible as a per-epoch $\Delta$ much larger than its training-only counterpart.

\paraq{Conclusion (frozen), with its risks.}
We freeze the following as this section's conclusion. It is stated as a disjunction so that it does
not depend on exactly where the pending cells land; the risks that bear on it are listed rather than
hidden, and any cell still awaiting measurement is marked in the tables and excluded from the
statement.

\begin{center}
\fbox{\begin{minipage}{0.95\linewidth}
\textbf{Conclusion.}\\[3pt]
Across both graphs, interconnects, and models, gsplit's per-epoch advantage over DGL is
\emph{not uniform}: where it exists it is modest and \emph{shrinks} as hidden width or depth grows,
and under NVLink it narrows toward the noise floor and in places \emph{reverses}. We claim no single
crossover point---only that an advantage which is absent or diminishing in exactly the
large-width\,/\,deep regimes of interest is not something to rely on, which is what motivates the
hybrid-parallelism search of Section~\ref{sec:hybrids}.\\[5pt]
\textbf{Risks to this conclusion.}
\begin{itemize}\itemsep1pt
  \item[\textbf{R1}] \emph{Protocol.} Any carried-over papers100M-NVLink rows predate the unified
        $3$-epoch\,/\,$1200$\,s protocol; magnitudes may shift on refresh (direction not expected to).
  \item[\textbf{R2}] \emph{Unmeasured corners.} The largest-$h$ / deepest cells may record
        \textsc{oom}/\textsc{t/o} rather than a number; either still supports ``not reliable.''
  \item[\textbf{R3}] \emph{Emulated PCIe.} PCIe is emulated via \texttt{NCCL\_P2P\_DISABLE=1}; its
        trends are indicative, not a native-PCIe measurement.
  \item[\textbf{R4}] \emph{What would overturn it.} Large and \emph{growing} gsplit gains across the
        large-$h$/deep cells; the measured trend to date runs against this.
\end{itemize}
The diminishing-gain and losing cells, once the matrix is complete, are enumerated in
\texttt{results/losing\_cells.csv} as the typed input to the simulator-driven configuration search of
Section~\ref{sec:simulator}.
\end{minipage}}
\end{center}
